\chapter{Introduction}

\section{Background}
Digital communication systems are widely used in modern monitoring and control applications 
where sensor data must be transmitted reliably from a sensing node to a receiver. 
In industrial environments, sensors are commonly deployed to monitor parameters 
such as temperature, \\vibration, pressure, gas concentration, flow rate, and machine health. 
These measurements are often converted into binary form and sent through a communication link for further processing, logging, or control action.

In this project, the selected application theme is 
\textbf{\textit{Industrial Sensor Monitoring \\Communication}}. 
In such a scenario, an industrial sensor node generates binary information and transmits it wirelessly to a nearby receiver. 
Because practical transmission channels are affected by \\attenuation, distortion, and noise, the received signal may differ from the transmitted signal. 
Hence, proper modulation and receiver design are essential for reliable data recovery.

Binary Phase Shift Keying (BPSK) is one of the simplest and most reliable digital modulation techniques for binary data transmission. 
In BPSK, bit `1` and bit `0` are represented by two antipodal carrier waveforms having a phase difference of \(180^\circ\). 
Due to its simple structure and strong noise tolerance, 
BPSK is widely used as a basic model for understanding digital communication theory.

The present work studies a complete BPSK communication chain using mathematical analysis and Python simulation. 
The work includes waveform generation, energy calculation, Fourier-domain analysis, bandwidth estimation, channel modeling, correlator detection, FIR matched filter implementation, BER simulation, and Shannon-capacity interpretation. 
Since assigned a mild low-pass channel followed by AWGN and both receiver types, this report gives special attention to channel distortion and detector comparison.

\section{Project Objective}
The main objective of this project is to develop a complete BPSK-based communication link in Python for the assigned industrial monitoring scenario 
and to analyse its performance using the mathematical concepts covered in the course.

The specific objectives of this project are listed below:
\begin{enumerate}[label=\arabic*.]
    \item To generate the BPSK symbols corresponding to binary `1` and binary `0` using the \\assigned carrier frequency, bit rate, and sampling frequency.
    \item To compute the bit duration, samples per bit, and energy per bit, and to compare numerical and theoretical energy values.
    \item To generate and visualize the transmitted 10-bit BPSK waveform for the prescribed test bit sequence.
    \item To analyse the FFT spectrum, power spectral density, bandwidth, and Parseval relation of the generated signal.
    \item To model the assigned transmission channel, namely a mild low-pass channel followed by additive white Gaussian noise.
    \item To implement and compare both a correlator receiver and an FIR matched filter receiver.
    \item To perform Monte Carlo BER simulation and compare simulated BER with the theoretical BER expression of BPSK.
    \item To compute Shannon channel capacity and discuss the gap between uncoded BPSK \\performance and the Shannon limit.
\end{enumerate}

\section{Course Outcome Coverage}

\begin{table}[H]
\centering
\caption{CO-wise project coverage}
\begin{tabular}{|c|p{0.75\linewidth}|}
\hline
\textbf{CO} & \textbf{Project Component} \\
\hline
CO1 & BPSK signal representation, sinusoidal modeling, antipodal signaling, and energy-per-bit calculation. \\
\hline
CO2 & Fourier spectrum analysis, power spectral density estimation, bandwidth calculation, and Parseval verification. \\
\hline
CO3 & Mild low-pass channel modeling, AWGN inclusion, convolution, and correlator receiver design. \\
\hline
CO4 & Sampled-signal processing, FIR matched filter design, stability discussion, and discrete-time interpretation using Z-transform ideas. \\
\hline
CO5 & Q-function based theoretical BER, random-bit generation, Monte Carlo BER simulation, and practical detector validation. \\
\hline
CO6 & Shannon capacity calculation, spectral-efficiency discussion, and comparison between practical uncoded BPSK and the Shannon limit. \\
\hline
\end{tabular}
\end{table}